\chapter{Revisão Bibliográfica}
\label{chapter:rev-bib}

% -------------------------------------------------------------
\section{Contextualização e Relevância da Previsão de Velocidade
         do Vento}
% -------------------------------------------------------------

A expansão das fontes de energia renovável nas últimas décadas
impõe desafios crescentes à operação e ao planejamento dos sistemas
elétricos de potência. A energia eólica, em particular, ocupa
posição de destaque nesse cenário, tendo registrado crescimento
expressivo de capacidade instalada globalmente ao longo dos anos
2000 e 2010, tendência que se acentuou na década de 2020
\cite{irena2024}. A natureza estocástica e intermitente da geração
eólica, contudo, representa um obstáculo fundamental à integração
dessas usinas à rede elétrica, uma vez que desequilíbrios entre
oferta e demanda comprometem a estabilidade do sistema, elevam
custos de despacho e aumentam a necessidade de reservas operativas.

Nesse contexto, a previsão acurada da velocidade do vento constitui
elemento crítico para a tomada de decisão em diferentes horizontes
temporais~--- do despacho intradiário à programação de manutenção
preventiva e ao planejamento de longo prazo. Erros de previsão
elevados resultam em custos de balanceamento significativos, despacho
subótimo de outras fontes e riscos à segurança operacional. Portanto,
o desenvolvimento de modelos de previsão de velocidade do vento com
maior precisão, interpretabilidade e robustez permanece como um
problema aberto de alta relevância científica e tecnológica.

A literatura científica sobre o tema tem crescido de forma notável.
Uma análise sistemática realizada no âmbito desta pesquisa,
compreendendo 893 artigos indexados na base Scopus e publicados
entre 2022 e 2025, evidencia trajetória ascendente de produção:
221~artigos em 2022, 193~em 2023, 232~em 2024 e projeção superior
a 240~para 2025. Esse volume de produção, concentrado em periódicos
de alto fator de impacto como \textit{Energy}, \textit{Renewable
Energy}, \textit{Energies} e \textit{Applied Energy}, reflete a
maturidade e a centralidade do tema na agenda de pesquisa em
engenharia de sistemas de energia.


% -------------------------------------------------------------
\section{Metodologia de Revisão Sistemática}
% -------------------------------------------------------------

\subsection{Estratégia de Busca e Constituição do Corpus}

A revisão sistemática foi conduzida na base de dados Scopus,
reconhecida pela abrangência e consistência de indexação em ciências
exatas e engenharias. A estratégia de busca adotou a seguinte
expressão de consulta:

\begin{lstlisting}[
    basicstyle=\small\ttfamily,
    breaklines=true,
    frame=single,
    caption={Expressão de busca utilizada na base Scopus.},
    label={lst:scopus-query}
]
TITLE-ABS-KEY (
  ( "wind speed forecasting" OR "wind speed prediction"
    OR "wind power forecasting" OR "wind power prediction" )
)
AND PUBYEAR > 2021
AND PUBYEAR < 2026
AND DOCTYPE ( ar )
AND ( LIMIT-TO ( LANGUAGE , "English" ) )
\end{lstlisting}

Os critérios de inclusão são: (i)~presença dos termos de busca no
título, resumo ou palavras-chave dos autores; (ii)~publicação entre
2022 e 2025; (iii)~tipologia documental restrita a artigos de
periódico (\textit{article}); e (iv)~idioma inglês. Optou-se por
restringir a busca a artigos originais~--- excluindo revisões,
anais de conferência, capítulos de livro e documentos de acesso
restrito~--- de modo a garantir que cada trabalho represente uma
contribuição metodológica inédita com processo de avaliação por
pares. A restrição temporal ao período 2022--2025 decorre da
estratégia deliberada de mapear o estado da arte recente, marcado
pela emergência e consolidação de arquiteturas Transformer e modelos
híbridos de decomposição--aprendizado profundo. O resultado bruto
da consulta, exportado no formato CSV pela interface de exportação
da Scopus, constituiu o insumo primário do processo de
classificação descrito na próxima seção.


% - - - - - - - - - - - - - - - - - - - - - - - - - - - - - -
\subsection{Ferramenta de Classificação Assistida por Inteligência
            Artificial: BIIA}

\subsubsection{Concepção e Arquitetura Geral}

A classificação temática de um corpus com centenas de artigos por
categorias multidimensionais~--- tarefa, variável-alvo, arquitetura
de rede neural, nome do modelo e técnicas complementares~--- é uma
tarefa que, se conduzida manualmente, demanda esforço considerável
e está sujeita à inconsistência intra-avaliador em julgamentos
limítrofes. Para enfrentar esse desafio, foi desenvolvida no âmbito
desta pesquisa a ferramenta \textbf{BIIA} (\textit{bibliografia +
inteligência artificial}), uma aplicação \textit{web} de página
única (\textit{single-page application}, SPA) implementada como
\textit{Progressive Web App} (PWA) e distribuída sob licença MIT.

O BIIA opera integralmente no navegador do usuário~--- sem
necessidade de instalação, servidor ou dependência de infraestrutura
remota~--- e pode ser utilizado \textit{offline} após o primeiro
carregamento, uma vez que seus recursos são armazenados em
\textit{cache} pelo \textit{service worker}. A interface, em tema
escuro, foi desenvolvida com ênfase na transparência do processo de
classificação: o usuário pode inspecionar, em tempo real, o
\textit{prompt} enviado à API para qualquer artigo antes de iniciar
o processamento, promovendo auditabilidade e reprodutibilidade do
protocolo.

A ferramenta recebe como entrada um arquivo CSV exportado
diretamente da Scopus e acrescenta ao arquivo original~---
preservando todas as colunas existentes~--- as colunas descritas
na Tabela~\ref{tab:biia-colunas}.

\begin{table}[htbp]
\centering
\caption{Colunas acrescentadas ao CSV de entrada pelo BIIA.}
\label{tab:biia-colunas}
\begin{tabular}{llp{7.5cm}}
\toprule
\textbf{Coluna} & \textbf{Tipo} & \textbf{Descrição} \\
\midrule
\texttt{task\_predicted}         & valor fechado   & Tarefa principal do artigo (11~categorias) \\
\texttt{target\_predicted}       & valor fechado   & Variável-alvo do modelo proposto (3~categorias) \\
\texttt{architecture}            & múltiplos rótulos & Famílias arquiteturais empregadas (15~categorias, combinável) \\
\texttt{model}                   & texto livre     & Nome ou acrônimo do modelo proposto \\
\texttt{complementary\_technique}& texto livre     & Técnicas de pré-processamento, otimização ou paradigma de treinamento \\
\texttt{task\_source}            & literal         & Origem da classificação de tarefa: \texttt{"regex"} ou \texttt{"ai"} \\
\texttt{target\_source}          & literal         & Origem da classificação de alvo: \texttt{"regex"} ou \texttt{"ai"} \\
\bottomrule
\end{tabular}
\end{table}


\subsubsection{Taxonomia de Classificação}

A taxonomia foi definida com base em revisão exploratória da
literatura e iterativamente refinada durante a fase piloto de
classificação, de modo a cobrir com precisão as linhas de pesquisa
presentes no corpus.

\paragraph{Tarefa (\textit{task}) — onze categorias mutuamente exclusivas:}

\begin{itemize}
  \item \texttt{previsão}: predição de velocidade do vento, potência
        ou variáveis meteorológicas em horizonte futuro definido;
  \item \texttt{otimização}: escalonamento, despacho, arbitragem de
        armazenamento, licitação em mercado de energia, seleção de
        atributos;
  \item \texttt{avaliação de recurso}: avaliação do potencial eólico
        de um sítio, ajuste de distribuição de Weibull, análise de
        viabilidade;
  \item \texttt{controle}: controle de esteira (\textit{wake}),
        controle de passo (\textit{pitch}), controle de frequência;
  \item \texttt{modelagem de curva}: curva de potência de turbina;
  \item \texttt{detecção de anomalias}: detecção de falhas, limpeza
        de dados SCADA;
  \item \texttt{geração de dados}: aumento de dados
        (\textit{data augmentation}), síntese de dados de vento;
  \item \texttt{simulação física}: modelo de fluidos computacional
        como contribuição principal (CFD, WRF, LES);
  \item \texttt{avaliação de risco}: quantificação de incerteza como
        objeto principal;
  \item \texttt{revisão}: revisão sistemática, \textit{survey},
        \textit{benchmark}, estudo comparativo;
  \item \texttt{outro}: artigos que não se enquadram nas categorias
        anteriores.
\end{itemize}

\paragraph{Alvo (\textit{target}) — três categorias hierarquicamente
ordenadas:}

\begin{enumerate}
  \item \texttt{velocidade do vento}: a variável de saída do modelo
        proposto é velocidade do vento;
  \item \texttt{potência}: a variável de saída é potência (de
        qualquer fonte) e não atende ao critério anterior;
  \item \texttt{outro}: qualquer outro caso.
\end{enumerate}

\paragraph{Arquitetura (\textit{architecture}) — quinze rótulos
atômicos e combináveis:}

Os rótulos de arquitetura refletem as famílias de modelos de
aprendizado profundo predominantes na literatura recente e podem
ser combinados para representar arquiteturas híbridas:

\begin{itemize}
  \item \texttt{Transformer} — Informer, PatchTST, Crossformer,
        Autoformer, Mamba;
  \item \texttt{Rede Recorrente} — LSTM, GRU, BiLSTM, xLSTM;
  \item \texttt{Rede Convolucional} — CNN, TCN, ResNet, ConvLSTM;
  \item \texttt{Rede de Grafos} — GCN, GAT, AGCRN, STGCN;
  \item \texttt{MLP-based} — N-BEATS, N-HiTS, TSMixer, TiDE;
  \item \texttt{Ensemble} — XGBoost, LightGBM, Random Forest,
        CatBoost;
  \item \texttt{Neuro-Fuzzy} — ANFIS;
  \item \texttt{LLM} — GPT, LLaMA aplicado a séries temporais;
  \item \texttt{Estatístico} — ARIMA, SARIMA, regressão, DLM
        Bayesiano;
  \item \texttt{Físico} — NWP, WRF, CFD, LES como componente de
        modelo híbrido;
  \item \texttt{Probabilístico} — modelos probabilísticos não
        enquadrados nas categorias acima;
  \item \texttt{Decomposição} — métodos baseados em decomposição
        como contribuição arquitetural;
  \item \texttt{Atenção} — mecanismo de atenção não integrado a
        Transformer;
  \item \texttt{Outro} — KAN, modelos de difusão, redes quânticas,
        aprendizado por reforço;
  \item \texttt{-} — ausência de modelo proposto identificável.
\end{itemize}

Os campos \texttt{model} e \texttt{complementary\_technique} são
preenchidos em texto livre. Metaheurísticas (PSO, GA, NSGA),
paradigmas de treinamento (aprendizado por transferência, aprendizado
federado, \textit{meta-learning}, destilação de conhecimento) e
estratégias metodológicas (regressão quantílica, Monte Carlo) são
classificados como \texttt{complementary\_technique}~--- não como
\texttt{architecture}~--- segundo regra explícita incorporada ao
\textit{prompt} de classificação.


\subsubsection{Pipeline de Classificação em Quatro Fases}

O processo de classificação é organizado em quatro fases sequenciais
que obedecem a uma lógica de custo-benefício: as fases
determinísticas, de custo computacional nulo, são executadas
primeiro; as fases probabilísticas, que consomem recursos de API,
operam apenas sobre os artigos não resolvidos anteriormente.

\paragraph{Fase~1 — Classificação por Expressões Regulares
(\textit{Regex}).}

A primeira fase aplica um conjunto de expressões regulares ao título
e às palavras-chave dos autores~--- excluindo deliberadamente o
resumo e os \textit{Index Keywords} da Scopus, que frequentemente
contêm menções contextuais, e não a contribuição principal do
artigo. Essa escolha de escopo conservador visa maximizar a
precisão, delegando casos ambíguos às fases seguintes.

As regras de tarefa são testadas em ordem de prioridade decrescente
de especificidade, conforme a Tabela~\ref{tab:regex-priority}. A
ordem é relevante: \texttt{otimização} é verificada antes de
\texttt{previsão} porque artigos de otimização frequentemente
mencionam previsão como subcomponente. O termo
\textit{``optimization''} isolado é excluído do gatilho de
\texttt{otimização} para evitar falsos positivos oriundos de artigos
que descrevem apenas a otimização de hiperparâmetros do modelo.
As regras de alvo seguem hierarquia equivalente: \texttt{velocidade
do vento} tem prioridade sobre \texttt{potência}. Artigos sem
correspondência em qualquer regra têm o campo correspondente
deixado em branco para processamento pelas fases seguintes.

\begin{table}[htbp]
\centering
\caption{Hierarquia de prioridade das regras de tarefa na Fase~1
         (Regex).}
\label{tab:regex-priority}
\begin{tabular}{clp{7cm}}
\toprule
\textbf{Prioridade} & \textbf{Categoria} & \textbf{Termos gatilho (título ou palavras-chave)} \\
\midrule
1 & \texttt{revisão}              & ``systematic review'', ``literature survey'' \\
2 & \texttt{simulação física}     & ``CFD'', ``LES'', ``WRF'' \\
3 & \texttt{detecção de anomalias}& ``anomaly detection'', ``fault detection'', ``fault diagnosis'' \\
4 & \texttt{geração de dados}     & ``data augmentation'', ``synthetic data'' \\
5 & \texttt{avaliação de risco}   & ``risk assessment'' \\
6 & \texttt{controle}             & ``wake control'', ``pitch control'', ``frequency control'' \\
7 & \texttt{modelagem de curva}   & ``power curve'' \\
8 & \texttt{avaliação de recurso} & ``resource assessment'', ``wind potential'' \\
9 & \texttt{otimização}           & ``scheduling'', ``dispatch'', ``arbitrage'', ``unit commitment'' \\
10& \texttt{previsão}             & ``forecasting'', ``prediction'' \\
\bottomrule
\end{tabular}
\end{table}

\paragraph{Fases~2, 3 e 4 — Classificação via API de Modelo de
Linguagem de Grande Porte.}

As três fases de API são executadas em paralelo por meio de
\texttt{Promise.all}, minimizando o tempo total de processamento.
Cada fase envia ao modelo apenas os blocos de \textit{prompt}
relevantes para as colunas que processa, reduzindo o consumo de
\textit{tokens}:

\begin{itemize}
  \item \textbf{Fase~2}: classifica \texttt{task\_predicted} para
        artigos sem resultado da Fase~1;
  \item \textbf{Fase~3}: classifica \texttt{target\_predicted} para
        artigos sem resultado da Fase~1;
  \item \textbf{Fase~4}: classifica \texttt{architecture},
        \texttt{model} e \texttt{complementary\_technique} para
        \emph{todos} os artigos (a Fase~1 não cobre essas
        dimensões).
\end{itemize}

O modelo de linguagem empregado é o \textbf{Google Gemini}
(\texttt{gemini-3.1-flash-lite-preview}), acessado via API REST
pública. A escolha se justifica pela disponibilidade na camada
gratuita (15~requisições por minuto), adequada ao volume do corpus,
e pelo desempenho demonstrado em tarefas de classificação de texto
científico com \textit{prompt} estruturado.

Os parâmetros de processamento em lote estão sumarizados na
Tabela~\ref{tab:batch-params}.

\begin{table}[htbp]
\centering
\caption{Parâmetros de processamento em lote nas Fases~2, 3 e~4.}
\label{tab:batch-params}
\begin{tabular}{ll}
\toprule
\textbf{Parâmetro} & \textbf{Valor} \\
\midrule
Artigos por requisição (\textit{batch size}) & 5 \\
Intervalo entre lotes                        & 6.000~ms \\
Temperatura                                  & 0 (saída determinística) \\
Máximo de \textit{tokens} na resposta        & 2.048 \\
\bottomrule
\end{tabular}
\end{table}

Em caso de falha de API (erro HTTP 429 — limite de taxa excedido —
ou erro de servidor), o sistema implementa protocolo de reenvio com
\textit{backoff} exponencial: aguarda 120~segundos antes das três
primeiras tentativas e 300~segundos a partir da quarta. O contador
de tentativas é exibido ao usuário com \textit{timer} de contagem
regressiva, e a classificação pode ser pausada a qualquer momento.

\paragraph{Engenharia de \textit{Prompt}.}

Cada \textit{prompt} é composto por blocos modulares correspondentes
às colunas selecionadas. A instrução de sistema apresenta o agente
como especialista em aprendizado de máquina aplicado à previsão de
séries temporais; em seguida, cada bloco de classificação define com
precisão a categoria a ser inferida, lista as opções permitidas com
exemplos representativos, enuncia as regras de desambiguação e
especifica o formato de saída. O formato de resposta exigido é JSON
estrito, com um objeto por artigo contendo a chave \texttt{index}
(posição no lote) e as colunas classificadas. O modelo é
explicitamente instruído a responder \textbf{apenas} com JSON
válido, sem texto adicional.

Notavelmente, o bloco de arquitetura inclui regra explícita
diferenciando arquitetura de técnica complementar~--- distinção
essencial para a qualidade taxonômica da classificação. Metaheurísticas
e paradigmas de treinamento são instruídos a ir para
\texttt{complementary\_technique}, e nunca para \texttt{architecture}.

Em caso de falha no \textit{parsing} estrito, um módulo de extração
por expressão regular tenta recuperar os objetos JSON individuais do
texto de resposta; se ambas as tentativas falharem, o lote é
ignorado com registro de aviso e pode ser reprocessado manualmente.

\paragraph{Modos de Preenchimento.}

O usuário pode selecionar, por coluna, um de dois modos de
preenchimento: \textit{preencher vazios} (padrão)~--- ignora
registros com valor não vazio, útil para retomar sessões
interrompidas~---; ou \textit{sobrescrever tudo}~--- reprocessa
todos os registros, útil após refinamento taxonômico.


\subsubsection{Fluxo de Operação da Interface}

O usuário opera a ferramenta em cinco etapas sequenciais:

\begin{enumerate}
  \item \textbf{Autenticação}: inserção da chave de API do Google
        Gemini (formato \texttt{AIza\ldots}, mínimo de 20
        caracteres), armazenada localmente no navegador
        (\textit{localStorage}) para sessões futuras;
  \item \textbf{Carregamento do CSV}: \textit{drag-and-drop} ou
        seleção via diálogo de arquivo; o arquivo é parseado
        integralmente no lado cliente, sem transmissão a servidores
        externos;
  \item \textbf{Configuração das colunas}: seleção das dimensões a
        classificar (individualmente habilitáveis), escolha do modo
        de preenchimento e inspeção prévia do \textit{prompt} e do
        \textit{payload} real do primeiro artigo;
  \item \textbf{Execução da classificação}: barras de progresso
        individuais por fase, contadores de artigos processados e
        \textit{timer} de reintento visíveis em tempo real; botão de
        pausa disponível a qualquer momento;
  \item \textbf{Exportação}: botão de \textit{download} gera o CSV
        resultante com todas as colunas originais preservadas e as
        novas colunas de classificação acrescentadas.
\end{enumerate}


% - - - - - - - - - - - - - - - - - - - - - - - - - - - - - -
\subsection{Revisão Manual e Controle de Qualidade}

A automatização do processo de classificação não dispensa a
verificação humana. Após a conclusão das quatro fases de
classificação automática, foi conduzida revisão manual de uma
amostra dos artigos classificados. Para cada artigo revisado, foram
adicionadas ao CSV as colunas \texttt{task\_checked} e
\texttt{target\_checked} (rótulo atribuído pelo especialista) e
\texttt{task\_isequal} e \texttt{target\_isequal} (indicadores
booleanos de concordância entre a classificação automática e a
revisão manual). Essa estrutura permite rastrear discordâncias de
forma sistemática e alimentar as análises de acurácia descritas na
seção seguinte.

A revisão abrangeu preferencialmente artigos classificados
exclusivamente pela Fase~1 (\textit{regex}), uma vez que os artigos
das fases de API tiveram cobertura de revisão integral pelo próprio
processo iterativo de desenvolvimento da ferramenta. Os critérios de
julgamento adotados foram os mesmos enunciados na taxonomia formal
(Seção~2.2.2), com decisão baseada na leitura do título,
palavras-chave e, quando necessário, do resumo completo.


% - - - - - - - - - - - - - - - - - - - - - - - - - - - - - -
\subsection{Análises Quantitativas e Estatísticas}

\subsubsection{Módulos de Análise}

Após a classificação e revisão, três módulos independentes de
análise em Python processam o CSV resultante, filtrando o corpus
para os artigos com \texttt{task} $=$ \texttt{"previsão"} e
\texttt{target} $=$ \texttt{"velocidade do vento"}~--- 893~artigos
no total.

\paragraph{Módulo Bibliométrico (\texttt{run\_bibliometric.py}).}
Produz os indicadores clássicos de análise bibliométrica:
distribuição temporal de publicações, frequência de palavras-chave
dos autores, rede de coocorrência de termos, nuvem de palavras dos
resumos e rankings de autores, periódicos, países e instituições
mais produtivos. Todos os resultados são exportados em CSV e PNG.

\paragraph{Módulo de Análise Estrutural (\texttt{run\_combined.py}).}
Concentra-se nas dimensões metodológicas extraídas pela classificação
automática. Gera: (i)~estatísticas gerais do corpus; (ii)~ranking
dos artigos mais citados por arquitetura; (iii)~relatório de lacunas
(\textit{gap signals}), identificando arquiteturas com dois ou menos
artigos e arquiteturas em ascensão nos últimos dois anos do período;
(iv)~matriz de coocorrência arquitetura\,$\times$\,técnica
complementar, visualizada como mapa de calor; (v)~evolução temporal
das arquiteturas; e (vi)~rede de coocorrência das trinta combinações
mais frequentes de técnicas complementares.

\paragraph{Módulo de Validação (\texttt{run\_predicted.py}).}
Avalia a qualidade do processo de classificação automática a partir
das colunas de revisão manual. Produz gráficos de cobertura,
acurácia por fonte (\textit{regex} vs.\ API), acurácia por categoria
de tarefa e de alvo, matrizes de confusão e análise detalhada de
discordâncias.


\subsubsection{Estimação de Acurácia com Intervalo de Confiança de
              Wilson}

Para artigos classificados exclusivamente pela Fase~1
(\textit{regex})~--- para os quais a revisão manual constitui uma
amostra, não um censo~--- a acurácia pontual é complementada pelo
\textbf{intervalo de confiança de Wilson} a 95\%, mais robusto que
o intervalo de Wald em amostras com acurácia próxima de 1
\cite{wilson1927}. Para uma amostra de tamanho $n$ com acurácia
observada $\hat{p}$ e $z = 1{,}96$, o intervalo é dado por:

\begin{equation}
\label{eq:wilson}
  IC_{95\%}
  =
  \left[
    \frac{
      \hat{p} + \dfrac{z^{2}}{2n}
      \;\pm\;
      z\,\sqrt{\dfrac{\hat{p}(1-\hat{p})}{n}
               + \dfrac{z^{2}}{4n^{2}}}
    }{
      1 + \dfrac{z^{2}}{n}
    }
  \right].
\end{equation}

O tamanho mínimo de amostra necessário para margem de erro de
$\pm$5~pontos percentuais a 95\% de confiança é calculado como:

\begin{equation}
\label{eq:nmin}
  n_{\min}
  =
  \left\lceil
    \frac{z^{2}\,\hat{p}(1-\hat{p})}{0{,}05^{2}}
  \right\rceil.
\end{equation}

Amostras com $n_{\text{revisados}} \geq n_{\min}$ são sinalizadas
como suficientes; as demais, como insuficientes.


\subsubsection{Teste de Representatividade ($\chi^{2}$)}

Para verificar se a amostra revisada (artigos \textit{regex}) é
representativa da população não revisada da mesma origem,
aplica-se o \textbf{teste qui-quadrado de independência} sobre a
tabela de contingência $2 \times k$~--- linhas: \textit{revisado}
vs.\ \textit{não revisado}; colunas: $k$ categorias preditas~---
com agregação de categorias cuja frequência esperada seja inferior
a~5. A hipótese nula é a de homogeneidade distribucional entre os
dois grupos:

\begin{equation}
\label{eq:chi2}
  \chi^{2}
  =
  \sum_{i=1}^{2}\sum_{j=1}^{k}
  \frac{(O_{ij} - E_{ij})^{2}}{E_{ij}},
\end{equation}

\noindent onde $O_{ij}$ e $E_{ij}$ são as frequências observada e
esperada na célula $(i,\,j)$, respectivamente. Valor-$p \geq 0{,}05$
é interpretado como indicativo de representatividade; valor-$p <
0{,}05$ sinaliza viés de seleção e deve ser considerado na
interpretação das estimativas de acurácia. Os resultados são
consolidados em tabela estratificada que reporta, para cada
combinação de estrato (origem da classificação) e dimensão (tarefa
ou alvo): cobertura percentual, acurácia pontual, intervalo de
Wilson (Equação~\ref{eq:wilson}), estatística $\chi^{2}$, graus de
liberdade, valor-$p$ e indicador de suficiência amostral.


\subsubsection{Síntese dos Produtos de Análise}

O processo descrito nesta seção produz mais de 35~arquivos de saída,
incluindo gráficos em PNG, tabelas em CSV, relatórios em Markdown e
matrizes de confusão. A Tabela~\ref{tab:outputs} apresenta uma
síntese dos principais produtos gerados por cada módulo.

\begin{longtable}{llp{6.5cm}}
\caption{Principais produtos de análise gerados pelos três módulos
         Python.}
\label{tab:outputs} \\
\toprule
\textbf{Módulo} & \textbf{Arquivo} & \textbf{Conteúdo} \\
\midrule
\endfirsthead
\multicolumn{3}{c}{\tablename~\thetable{} — continuação} \\
\toprule
\textbf{Módulo} & \textbf{Arquivo} & \textbf{Conteúdo} \\
\midrule
\endhead
\midrule
\multicolumn{3}{r}{\textit{continua na próxima página}} \\
\endfoot
\bottomrule
\endlastfoot
% Bibliométrico
\multirow{7}{*}{\texttt{bibliometric}}
 & \texttt{papers\_per\_year.png/.csv}     & Distribuição temporal de publicações \\
 & \texttt{keywords.png/.csv}              & Frequência das palavras-chave dos autores \\
 & \texttt{keywords\_cooccurrence.png}     & Rede de coocorrência de termos \\
 & \texttt{wordcloud.png}                  & Nuvem de palavras dos resumos \\
 & \texttt{top\_authors.csv}               & Autores mais produtivos \\
 & \texttt{top\_journals.csv}              & Periódicos mais frequentes \\
 & \texttt{top\_countries.csv}             & Países de afiliação dos autores \\
\midrule
% Combined
\multirow{7}{*}{\texttt{combined}}
 & \texttt{overview.csv}                   & Estatísticas gerais do corpus \\
 & \texttt{gap\_signals.md}                & Arquiteturas raras e emergentes \\
 & \texttt{arch\_x\_complementary.png/.csv}& Mapa de calor arquitetura\,$\times$\,técnica \\
 & \texttt{arch\_trend.png/.csv}           & Evolução temporal das arquiteturas \\
 & \texttt{top\_cited.csv}                 & Artigos mais citados por arquitetura \\
 & \texttt{complementary\_cooccurrence\_network.png} & Rede de coocorrência de técnicas \\
 & \texttt{complementary\_cooccurrence\_matrix.png}  & Matriz de coocorrência de técnicas \\
\midrule
% Predicted
\multirow{8}{*}{\texttt{predicted}}
 & \texttt{coverage.png}                   & Cobertura da revisão manual por categoria \\
 & \texttt{accuracy\_by\_source.png}       & Acurácia por origem (\textit{regex} vs.\ API) \\
 & \texttt{accuracy\_by\_task\_category.png}  & Acurácia por categoria de tarefa \\
 & \texttt{accuracy\_by\_target\_category.png}& Acurácia por categoria de alvo \\
 & \texttt{confusion\_task.png/.csv}       & Matriz de confusão (tarefa) \\
 & \texttt{confusion\_target.png/.csv}     & Matriz de confusão (alvo) \\
 & \texttt{discrepancies.png}              & Principais discordâncias entre classificação e revisão \\
 & \texttt{statistical\_relevance\_stratified.csv} & Intervalo de Wilson e teste $\chi^{2}$ estratificados \\
\end{longtable}


% -------------------------------------------------------------
% (Seções 2.3–2.7 continuam conforme texto anterior)
% -------------------------------------------------------------

% =============================================================
% Referências utilizadas neste capítulo (adicionar ao .bib):
%
% @misc{irena2024,
%   author = {{IRENA}},
%   title  = {Renewable Power Generation Costs in 2023},
%   year   = {2024},
%   note   = {International Renewable Energy Agency}
% }
%
% @article{wilson1927,
%   author  = {Wilson, Edwin B.},
%   title   = {Probable inference, the law of succession, and
%               statistical inference},
%   journal = {Journal of the American Statistical Association},
%   volume  = {22},
%   number  = {158},
%   pages   = {209--212},
%   year    = {1927}
% }
% =============================================================
